\documentclass[11pt, a4paper]{article}

\usepackage[T1]{fontenc}
\usepackage[utf8]{inputenc}
\usepackage{tgpagella}
\usepackage[spanish, es-noshorthands]{babel}
\usepackage{amsmath, amssymb, bm}
\usepackage{booktabs}
\usepackage{tabularx}
\usepackage[table, xcdraw]{xcolor}
\usepackage[most]{tcolorbox}
\usepackage[american]{circuitikz}
\usepackage[margin=2.5cm]{geometry}
\usepackage{fancyhdr}

\pagestyle{fancy}
\fancyhf{}
\setlength{\headheight}{16pt}
\lhead{\textbf{UCB Sede Tarija} | Ing. Mecatrónica}
\chead{}
\rhead{IMT-121 Circuitos Electrónicos I | Prelaboratorio}
\lfoot{Prof: M.Sc. Ing. Bernardo Quiroga Turdera}
\rfoot{Página \thepage}
\renewcommand{\headrulewidth}{0.4pt}

\definecolor{ucbblue}{RGB}{0, 51, 102}

\begin{document}

\begin{tcolorbox}[colback=ucbblue!5!white, colframe=ucbblue, title=\textbf{Experiencia Integrada A: Prelaboratorio Obligatorio}]
    \centering \textbf{Cadena de Análisis Analítico, Simulación y Preparación para Medición}
\end{tcolorbox}

\textbf{Circuito Base:}
\begin{center}
\begin{circuitikz}[scale=1, transform shape, thick]
    \draw (0,0) to[V, l={$V_s=12\,\text{V}$}] (0,2) to[R, l={$R_1=4\,\text{k}\Omega$}] (3,2);
    \draw (3,2) to[short, *-o] (5,2) node[above]{$a$};
    \draw (5,2) to[R, l={$R_2=8\,\text{k}\Omega$}] (5,0);
    \draw (3,0) to[I, l={$I_s=1\,\text{mA}$}] (3,2);
    \draw (0,0) to[short, -o] (5,0) node[below]{$b$};
    \draw (2.5,0) node[ground]{};
\end{circuitikz}
\end{center}

\textbf{Requisito de Acceso:} Cada equipo debe presentar este documento debidamente calculado e implementado en simulación \textbf{antes} de ser autorizado a montar el circuito en protoboard.

\section*{Parte 1: Análisis Teórico de la Red}
\begin{enumerate}
    \item \textbf{Superposición para $\bm{V_{oc}}$:} Calcule la tensión de circuito abierto entre $a-b$ analizando el efecto de $V_s$ e $I_s$ por separado. Especifique las dos contribuciones y la suma algebraica ($V_{th}$).
    \vspace{2.5cm}
    \item \textbf{Cálculo de $\bm{R_{th}}$:} Redibuje la red con las fuentes independientes desactivadas y calcule la resistencia equivalente vista desde el puerto $a-b$.
    \vspace{2.5cm}
    \item \textbf{Predicción Norton:} Utilizando transformaciones de fuentes o relación directa, determine analíticamente la corriente teórica de cortocircuito $\bm{I_{sc}} = I_N$.
    \vspace{1.5cm}
    \item \textbf{Respuesta en Carga:} Si se conecta una carga de $R_L = 8\,\text{k}\Omega$, prediga $V_L$ e $I_L$.
    \vspace{1.5cm}
\end{enumerate}

\section*{Parte 2: Simulación SPICE y Seguridad}
\begin{enumerate}
    \item Adjunte capturas de pantalla de la simulación (.op) que verifiquen el voltaje de circuito abierto y la respuesta con la carga conectada.
    \item Escriba el procedimiento mediante el cual barrerá el valor de $R_L$ (usando comandos como \texttt{.step param}) para graficar $P_L$ vs $R_L$ en el simulador.
\end{enumerate}

\begin{tcolorbox}[colback=red!5!white, colframe=red!80!black, title=Gate de Seguridad Experimental para $\bm{I_{sc}}$]
    \textbf{Alto:} Para el montaje físico, se medirá $V_{oc}$ y la respuesta a varias cargas $R_L$. Sin embargo, la medición directa de cortocircuito con amperímetro en los terminales $a-b$ dependerá \textbf{exclusivamente} de la revisión del instructor del equipo real en mesa. Prepare su tabla de datos, pero marque $I_{sc\_medido}$ como "Pendiente Autorización".
\end{tcolorbox}

\end{document}
